%% ============================================================
%%  Chapitre 2 — État de l'art
%% ============================================================
\chapter{État de l'art}
\label{ch:etat_art}

\section{Grands modèles de langage et agents LLM}

\subsection{Des LLM aux agents}

Les grands modèles de langage (LLM), tels que GPT-4 \cite{openai2023gpt4},
Claude \cite{anthropic2024claude} ou LLaMA \cite{touvron2023llama}, ont
transformé le traitement automatique du langage naturel. Ces modèles, pré-entraînés
sur de vastes corpus textuels, exhibent des capacités remarquables en matière de
compréhension, de génération et de raisonnement. Cependant, leur utilisation
isolée présente des limitations importantes : absence de mémoire persistante
entre les appels, incapacité à accéder à des sources d'information externes
en temps réel, et difficulté à décomposer des tâches complexes en sous-tâches
gérables.

La notion d'\textbf{agent LLM} adresse ces limitations en dotant le LLM de
capacités supplémentaires : accès à des outils (\textit{tool use}),
mémoire externe, et capacité à planifier une séquence d'actions. Un agent
est formellement défini comme un système capable de percevoir son environnement
et d'agir sur lui pour atteindre un objectif \cite{wooldridge1995intelligent}.
Dans le contexte des LLM, un agent correspond à un modèle de langage augmenté
d'un ensemble d'outils (API, bases de données, moteurs de recherche) et d'un
mécanisme de contrôle déterminant quand et comment les utiliser.

\subsection{Patterns de raisonnement des agents}

Plusieurs paradigmes de raisonnement ont été proposés pour structurer le
comportement des agents LLM.

\paragraph{ReAct (Reason + Act).}
Yao et al. \cite{yao2022react} proposent un cadre interleaved où l'agent
alterne entre des étapes de \textit{Thought} (raisonnement verbalisé),
d'\textit{Action} (appel à un outil) et d'\textit{Observation} (résultat
de l'outil). Ce pattern améliore significativement la précision sur des
tâches de questions-réponses nécessitant des recherches factuelles,
tout en rendant le processus de raisonnement interprétable.

\paragraph{Reflexion.}
Shinn et al. \cite{shinn2023reflexion} introduisent Reflexion, un mécanisme
par lequel un agent LLM peut s'auto-évaluer et réviser ses productions en
s'appuyant sur un signal de feedback verbal. Contrairement au fine-tuning
traditionnel qui modifie les poids du modèle, Reflexion stocke les retours
dans la mémoire épisodique et les exploite lors des itérations suivantes.
Ce pattern présente l'avantage d'améliorer la qualité des sorties sans
nécessiter d'entraînement supplémentaire.

\paragraph{Tree-of-Thought (ToT).}
Yao et al. \cite{yao2023tree} généralisent la chaîne de pensée (\textit{chain
of thought}) à une exploration arborescente de l'espace de raisonnement. L'agent
génère plusieurs chemins de pensée en parallèle, les évalue, et sélectionne le
plus prometteur. Ce paradigme excelle pour les problèmes de planification à
plusieurs étapes.

\paragraph{Self-Consistency.}
Wang et al. \cite{wang2022self} montrent que la génération de plusieurs
chemins de raisonnement indépendants suivie d'un vote majoritaire améliore
la précision des LLM sur des tâches de raisonnement arithmétique et logique.
Ce principe est à la base des approches de Best-of-N que nous envisageons
pour les travaux futurs.

\subsection{Systèmes multi-agents LLM}

Au-delà des agents individuels, les \textbf{systèmes multi-agents LLM}
(LLM-MAS) mettent en jeu plusieurs agents spécialisés qui collaborent
pour résoudre des tâches complexes. Wang et al. \cite{wang2024survey}
proposent une taxonomie distinguant les architectures coopératives
(les agents poursuivent un objectif commun) des architectures compétitives
(les agents débattent pour converger vers une réponse de meilleure qualité).

Du et al. \cite{du2023improving} démontrent que le débat multi-agents
améliore la précision factuelle des LLM sur des benchmarks de mathématiques
et de raisonnement. Les agents génèrent indépendamment leurs réponses,
puis les échangent itérativement pour raffiner leurs positions jusqu'à
convergence.

\section{Orchestration d'agents : LangGraph}

\subsection{LangChain et l'écosystème agent}

LangChain \cite{chase2022langchain} est un framework Python qui simplifie
le chaînage d'appels aux LLM et l'intégration d'outils. Il propose des
abstractions pour les chaînes (\textit{chains}), les mémoires (\textit{memory})
et les agents. Cependant, pour des pipelines complexes avec des boucles,
des conditions et un état partagé entre agents, LangChain seul s'avère
insuffisant.

\subsection{LangGraph : graphes d'état pour les agents}

LangGraph \cite{langgraph2024} est une extension de LangChain qui modélise
les pipelines agents comme des graphes orientés dont les nœuds sont des
agents et les arêtes représentent les transitions d'état. L'état du système
est un dictionnaire partagé (\texttt{TypedDict}) que chaque agent peut lire
et écrire, dans le respect de la règle \textit{Single Writer / Multiple
Readers}. Cette architecture offre trois avantages majeurs par rapport
aux approches séquentielles : la gestion naturelle des boucles (Reflexion,
ReAct), la possibilité d'arêtes conditionnelles (\textit{conditional edges})
pour router dynamiquement le flux selon l'état, et la persistance de l'état
entre les nœuds sans besoin de variables globales.

\section{Qualité des publications scientifiques}

\subsection{Indicateurs bibliométriques}

L'évaluation de la qualité d'une publication académique s'appuie
traditionnellement sur plusieurs indicateurs bibliométriques.

\paragraph{Impact de la revue.}
Le facteur d'impact (\textit{Impact Factor}) \cite{garfield1972citation}
et le classement par quartile (Q1, Q2, Q3, Q4 selon Scimago ou Web of Science)
mesurent la visibilité d'une revue dans sa discipline. Un article publié dans
une revue de rang Q1 bénéficie d'une présomption de qualité supérieure,
bien que ce critère doive être nuancé par la culture de publication propre
à chaque discipline.

\paragraph{Citations.}
Le nombre de citations d'un article constitue le proxy le plus commun de
son impact \cite{hirsch2005index}. L'indice H (\textit{h-index}) de Hirsch,
qui mesure la productivité et l'influence d'un chercheur, est un indicateur
complémentaire permettant d'évaluer l'autorité scientifique des auteurs
d'un article.

\paragraph{Vélocité des citations.}
Au-delà du compte brut, la \textit{vélocité de citation} — nombre de citations
récentes par unité de temps — permet d'identifier les travaux émergents à fort
impact potentiel, notamment les articles récents qui accumulent rapidement
des citations \cite{ke2015defining}.

\subsection{Pertinence sémantique}

La pertinence sémantique d'un article par rapport à un sujet de recherche
peut être mesurée à deux niveaux. L'approche lexicale, basée sur les
n-grammes et la correspondance de mots-clés \cite{jones1972statistical},
est rapide mais souffre du problème de synonymie et de polysémie.
L'approche par embeddings, fondée sur des modèles de langage comme
\textit{Sentence-BERT} \cite{reimers2019sentencebert} ou ses variantes
légères (all-MiniLM-L6-v2 \cite{wang2020minilm}), capture la similarité
sémantique dans un espace vectoriel continu et s'avère bien plus robuste.

\section{AutoML et optimisation d'hyperparamètres}

\subsection{Optimisation bayésienne}

L'optimisation bayésienne \cite{snoek2012practical} est une méthode de
recherche d'hyperparamètres fondée sur la construction d'un modèle de
substitution (\textit{surrogate model}) de la fonction objectif. À chaque
itération, l'algorithme sélectionne les hyperparamètres les plus prometteurs
en équilibrant exploration (zones peu connues de l'espace) et exploitation
(zones connues pour leur haute performance), via une fonction d'acquisition
telle que l'espérance d'amélioration (\textit{Expected Improvement}).

\subsection{Optuna et l'échantillonneur TPE}

Optuna \cite{akiba2019optuna} est un framework AutoML Python qui implémente
l'algorithme TPE (\textit{Tree-structured Parzen Estimator}) de Bergstra et
al.\ \cite{bergstra2011algorithms}. TPE modélise séparément la distribution
des hyperparamètres associés aux bons résultats $l(x)$ et aux mauvais résultats
$g(x)$, puis sélectionne les candidats qui maximisent le ratio $l(x)/g(x)$.
Comparé à la recherche aléatoire ou par grille, TPE converge vers de bonnes
solutions avec significativement moins d'évaluations de la fonction objectif,
ce qui est crucial lorsque chaque évaluation nécessite un appel API coûteux.

\section{Méta-learning}

\subsection{Définition et taxonomie}

Le méta-learning, ou \textit{apprentissage à apprendre}, désigne un ensemble
de techniques visant à améliorer les performances d'un algorithme d'apprentissage
en exploitant l'expérience accumulée sur de multiples tâches passées
\cite{vanschoren2018meta}. Contrairement à l'apprentissage conventionnel qui
optimise les performances sur une tâche spécifique, le méta-learning cherche à
apprendre une représentation ou une procédure générale transférable à de
nouvelles tâches avec peu d'exemples.

Trois grandes familles d'approches existent :
\begin{itemize}
  \item Les méthodes fondées sur l'optimisation du modèle initial (MAML
    \cite{finn2017model} et ses variantes), qui apprennent une initialisation
    des poids du modèle facilitant le fine-tuning rapide sur de nouvelles tâches ;
  \item Les méthodes fondées sur des métriques (Prototypical Networks
    \cite{snell2017prototypical}), qui apprennent un espace d'embedding dans
    lequel la classification par plus proche voisin est efficace ;
  \item Les méthodes fondées sur des modèles (\textit{model-based}), qui apprennent
    un mécanisme d'inférence rapide, souvent sous la forme d'un réseau de
    mémoire ou d'un transformeur.
\end{itemize}

\subsection{Méta-learning pour la configuration de systèmes}

Notre travail s'inscrit dans un paradigme de méta-learning différent des
approches précédentes, plus proche du domaine AutoML \cite{hutter2019automated}
et du \textit{meta-learning for algorithm configuration}
\cite{feurer2015initializing}. L'idée est de prédire les hyperparamètres
optimaux d'un algorithme (ici, les pondérations du scorer) pour une nouvelle
tâche (ici, un nouveau domaine de recherche), à partir de caractéristiques
descriptives de cette tâche (\textit{meta-features}).

Feurer et al.\ \cite{feurer2015initializing} utilisent les meta-features
d'un dataset (nombre d'instances, de features, ratio de classes, etc.) pour
initialiser intelligemment la recherche d'hyperparamètres dans Auto-sklearn.
Dans notre contexte, les meta-features décrivent le corpus de publications
d'un domaine (distribution des citations, niveau des auteurs, récence des
articles, etc.) et permettent de prédire les pondérations optimales du scorer.

\subsection{Combinaison d'estimateurs : agrégation de rangs}

La combinaison de plusieurs estimateurs (\textit{ensemble methods}) améliore
typiquement la robustesse et la performance par rapport à un estimateur unique.
Dans le contexte de la prédiction de vecteurs de pondérations, nous adaptons
des méthodes d'agrégation de rangs initialement développées pour la fusion
de résultats de recherche d'information.

Le \textbf{comptage de Borda} \cite{borda1784memoire}, formalisé dans un
contexte d'apprentissage par Dwork et al.\ \cite{dwork2001rank}, attribue
à chaque alternative un score égal à sa position dans le classement inverse,
puis somme ces scores sur tous les votants. Il s'agit d'une méthode robuste
aux outliers et naturellement insensible aux différences d'échelle entre
estimateurs.

La \textbf{fusion par rang réciproque} (RRF, \textit{Reciprocal Rank Fusion})
\cite{cormack2009reciprocal} assigne à chaque alternative un score
$\text{RRF}(d) = \sum_{r \in R} \frac{1}{k + \text{rank}_r(d)}$ où $k$ est
un paramètre de lissage. RRF est particulièrement efficace pour la fusion
de listes de rangs provenant de systèmes hétérogènes.

\section{Métriques d'évaluation en recherche d'information}

\subsection{NDCG (\textit{Normalized Discounted Cumulative Gain})}

La métrique NDCG \cite{jarvelin2002cumulated} est le standard d'évaluation
des systèmes de classement en recherche d'information. Elle prend en compte
à la fois la pertinence et la position des documents dans le classement,
en appliquant un facteur de discount logarithmique aux positions inférieures.

Pour une liste de $k$ documents, le DCG est défini par :
\begin{equation}
  \text{DCG}@k = \sum_{i=1}^{k} \frac{2^{r_i} - 1}{\log_2(i+1)}
  \label{eq:dcg}
\end{equation}
où $r_i \in \{0, 1, 2\}$ est la pertinence du document en position $i$
(0 = non pertinent, 1 = pertinent, 2 = très pertinent).

Le NDCG est obtenu en normalisant le DCG par le DCG idéal (IDCG) :
\begin{equation}
  \text{NDCG}@k = \frac{\text{DCG}@k}{\text{IDCG}@k}
  \label{eq:ndcg}
\end{equation}

Nous utilisons NDCG@15 comme métrique principale, ce choix étant motivé
par la pratique usuelle de consulter les 15 à 20 premiers résultats d'une
recherche académique.

\subsection{Tests statistiques non-paramétriques}
\label{sec:etat_art_wilcoxon}

Le \textbf{test de Wilcoxon signé} \cite{wilcoxon1945individual} est une
alternative non-paramétrique au test $t$ apparié, adaptée aux situations
où les hypothèses de normalité ne peuvent être vérifiées. Pour $n$ paires
d'observations $(x_i, y_i)$, on considère les différences $d_i = y_i - x_i$,
on les classe par valeur absolue, et on calcule la statistique $W^+$ (somme
des rangs des différences positives). La corrélation rang-bisériale
$r_{rb} = (W^+ - W^-) / (W^+ + W^-)$ mesure la taille d'effet, avec les
conventions : $|r_{rb}| < 0.3$ (petit), $0.3$--$0.5$ (moyen), $> 0.5$ (grand).

\section{Optimisation automatique des prompts}

\subsection{L'importance du \textit{prompt engineering}}

La qualité des sorties d'un LLM dépend fortement de la formulation des
instructions (\textit{prompts}) qui lui sont soumises. Le \textit{prompt
engineering} manuel est un processus itératif et subjectif, difficile à
systématiser. Plusieurs travaux ont proposé d'automatiser cette optimisation.

\subsection{ProTeGi}

Pryzant et al.\ \cite{pryzant2023automatic} proposent ProTeGi
(\textit{Prompt Optimization with Textual Gradients}), une méthode qui
adapte le concept de gradient de descente au domaine textuel. ProTeGi
utilise un LLM ``critique'' pour identifier les faiblesses d'un prompt
et proposer des modifications ciblées, de façon analogue à un gradient
qui indique la direction d'amélioration. Cette approche ne nécessite
pas d'accès aux gradients du modèle et est donc applicable à tout LLM
accessible via API.

\section{Veille scientifique automatisée — travaux existants}

\subsection{Systèmes de recommandation académique}

Plusieurs systèmes automatisés de veille ou de recommandation académique
ont été proposés. \textit{Semantic Scholar} \cite{lo2020s2orc} utilise des
embeddings pour recommander des articles similaires. \textit{ResearchRabbit}
et \textit{Connected Papers} proposent des visualisations de graphes de
citations pour explorer des corpus. Ces systèmes se concentrent sur la
découverte d'articles similaires à un corpus de départ, sans aborder
l'analyse automatique du contenu ni la synthèse.

\subsection{Génération automatique de revues de littérature}

Des travaux récents ont exploré la génération automatique de revues de
littérature (\textit{automated survey generation}) à l'aide de LLM.
Liu et al.\ \cite{liu2023survey} utilisent un agent LLM pour décomposer
un sujet de recherche en sous-thèmes et récupérer les articles pertinents
pour chacun. Cependant, ces approches ne proposent pas de mécanisme de
filtrage par qualité ni d'évaluation rigoureuse de la généralisation à de
nouveaux domaines.

\subsection{Positionnement de notre travail}

Par rapport aux travaux existants, notre contribution se distingue sur
trois axes :
\begin{enumerate}
  \item \textbf{Scoring multi-dimensionnel et calibration par AutoML} : nous
    formalisons l'évaluation de la qualité d'un article selon six dimensions
    et apprenons automatiquement les pondérations optimales par domaine.
  \item \textbf{Architecture multi-agents complète} : nous proposons un pipeline
    end-to-end intégrant sept agents spécialisés, des patterns de raisonnement
    avancés (ReAct, Reflexion) et un contrat de communication formalisé.
  \item \textbf{Méta-learning pour la généralisation} : nous évaluons
    rigoureusement, sur 19 domaines, la capacité du système à s'adapter
    à de nouveaux domaines sans reconfiguration manuelle.
\end{enumerate}

Le tableau~\ref{tab:related_work} synthétise le positionnement de notre
approche par rapport aux systèmes existants.

\begin{table}[htbp]
  \centering
  \caption{Comparaison des approches de veille scientifique automatisée.}
  \label{tab:related_work}
  \rowcolors{2}{lightgray}{white}
  \begin{tabular}{lcccc}
    \toprule
    \textbf{Système} & \textbf{Multi-agents} & \textbf{Scoring} & \textbf{Synthèse} & \textbf{Généralisation} \\
    \midrule
    Semantic Scholar & \texttimes & partiel & \texttimes & \checkmark \\
    Liu et al. (2023) & \checkmark & \texttimes & \checkmark & \texttimes \\
    ResearchRabbit & \texttimes & \texttimes & \texttimes & \checkmark \\
    \textbf{Notre approche} & \checkmark & \checkmark & \checkmark & \checkmark \\
    \bottomrule
  \end{tabular}
\end{table}
